\chapter{Conception}
\label{ch:conception}
\minitoc
\clearpage

\section{Introduction}

Concevoir PharmacieConnect a d'abord consisté à répartir les besoins fonctionnels exposés au chapitre
précédent entre les différentes briques de la solution. Le backend FastAPI assume les règles métier, la
sécurité, les imports et l'accès aux données. L'interface React/Vite prend en charge les opérations
d'administration, pendant que l'application Expo/React Native reste dédiée au parcours citoyen. Ce découpage
a été choisi pour ne pas dupliquer les règles sensibles dans les clients et pour conserver une architecture
plus simple à tester et à faire évoluer.

\section{Architecture globale}

L'architecture de PharmacieConnect place FastAPI comme point de passage unique entre les interfaces et les
données. Le tableau de bord React/Vite utilise cette API pour les opérations de gestion, alors que
l'application Expo/React Native l'utilise principalement pour la recherche, les gardes, les médicaments et le
compte utilisateur. Cette organisation évite que les règles de validation, d'audit et de restriction régionale
soient dupliquées dans les clients, comme l'illustre la Figure~\ref{fig:architecture-globale}.

\begin{figure}[H]
    \centering
    \includegraphics[width=0.95\textwidth]{plantuml/architecture_globale}
    \caption{Architecture globale de PharmacieConnect\\ \small Source : Auteur}
    \label{fig:architecture-globale}
\end{figure}

L'application mobile utilise principalement les endpoints publics de consultation, les routes liées au
compte utilisateur et l'API de l'assistant. Le panneau d'administration s'appuie sur des endpoints
protégés pour les opérations de gestion : ajout, mise à jour et retrait d'entrées, import CSV, consultation
des indicateurs et audit. Le backend devient ainsi la couche commune qui applique les règles métier avant toute
lecture ou écriture persistante.

\section{Backend FastAPI}

Dans le projet, FastAPI n'a pas été utilisé uniquement pour exposer des routes REST \cite{fastapi}. Il
constitue la couche qui relie les modules essentiels de PharmacieConnect : authentification, pharmacies,
gardes, médicaments, imports CSV, audit, indicateurs et assistant de premiers secours. Les routes restent
séparées des services métier afin que les règles de validation, les restrictions régionales et les accès à
la base via SQLAlchemy soient centralisés dans le backend plutôt que dispersés entre les interfaces. Le
Tableau~\ref{tab:backend-organisation} présente cette organisation.

\begin{table}[H]
    \centering
    \caption{Organisation du backend\\ \small Source : Auteur}
    \label{tab:backend-organisation}
    \begin{tabular}{L{4cm}L{9cm}}
        \toprule
        \textbf{Élément} & \textbf{Rôle} \\
        \midrule
        \texttt{main.py} & Initialisation de l'application FastAPI, middlewares, migrations au démarrage et endpoints publics. \\
        \texttt{routers/} & Routes d'authentification, d'administration et d'indicateurs. \\
        \texttt{services/} & Logique métier liée à l'authentification, aux pharmacies, aux gardes, aux médicaments, à l'audit et au cache. \\
        \texttt{models.py} & Modèles SQLAlchemy représentant les tables principales. \\
        \texttt{schemas.py} & Schémas Pydantic utilisés pour valider les entrées et structurer les sorties. \\
        \texttt{migrations/} & Scripts SQL permettant de faire évoluer le schéma de la base. \\
        \texttt{events/} & Bus d'événements interne et écouteurs associés. \\
        \bottomrule
    \end{tabular}
\end{table}

Le traitement d'une requête suit généralement le même parcours. Le client appelle une route FastAPI ; la
route vérifie les paramètres et l'utilisateur courant ; le service métier applique ensuite les règles
nécessaires avant l'accès à la base via SQLAlchemy.

\section{Tableau de bord React/Vite}

Côté administration, React et Vite ont été retenus pour construire une interface que l'on peut faire évoluer
page par page \cite{react,vite}. Dans notre cas, cette modularité est utile parce que le tableau de bord ne
sert pas seulement à afficher des informations : il doit aussi permettre de corriger les registres,
d'importer des fichiers, de vérifier des indicateurs et de suivre les actions sensibles (voir
Tableau~\ref{tab:admin-organisation}).

\begin{table}[H]
    \centering
    \caption{Organisation de l'administration web\\ \small Source : Auteur}
    \label{tab:admin-organisation}
    \begin{tabular}{L{4cm}L{9cm}}
        \toprule
        \textbf{Élément} & \textbf{Rôle} \\
        \midrule
        \texttt{src/pages/} & Pages principales du tableau de bord : pharmacies, gardes, médicaments, carte, audit, paramètres et gestion. \\
        \texttt{src/components/} & Composants d'interface, tableaux, formulaires, modales, cartes statistiques et éléments de layout. \\
        \texttt{src/context/} & Gestion de l'authentification, de la langue et de certains états globaux. \\
        \texttt{src/lib/api.js} & Client Axios responsable des appels backend et de l'injection du jeton JWT. \\
        \texttt{src/styles/} & Styles et utilitaires liés à la présentation de l'interface. \\
        \bottomrule
    \end{tabular}
\end{table}

Cette organisation évite de mélanger l'affichage avec la logique d'accès aux données. Le client API ajoute
le jeton d'accès aux requêtes protégées et prend en charge le renouvellement de session lorsque le jeton
expire.

\section{Application mobile Expo/React Native}

Pour la partie mobile, Expo et React Native ont été choisis afin de développer rapidement des écrans centrés
sur le parcours citoyen \cite{expo,reactnative}. L'organisation retenue suit les actions que l'utilisateur
réalise réellement dans l'application : chercher une pharmacie, ouvrir la carte, choisir une date de garde,
consulter un médicament, gérer son compte ou accéder à l'assistant de premiers secours, comme indiqué dans
le Tableau~\ref{tab:mobile-organisation}.

\begin{table}[H]
    \centering
    \caption{Organisation de l'application mobile\\ \small Source : Auteur}
    \label{tab:mobile-organisation}
    \begin{tabular}{L{4cm}L{9cm}}
        \toprule
        \textbf{Élément} & \textbf{Rôle} \\
        \midrule
        \texttt{screens/} & Écrans de l'application : accueil, carte, calendrier, médicaments, connexion, inscription, vérification email et paramètres. \\
        \texttt{components/} & Composants réutilisables comme les cartes de pharmacies, modales, sélecteurs et éléments RTL. \\
        \texttt{config/api.js} & Configuration de l'URL backend et des appels API. \\
        \texttt{utils/} & Fonctions utilitaires pour les données, les couleurs, le thème, le logging et la gestion RTL. \\
        \texttt{locales/} & Fichiers de traduction français, anglais et arabe. \\
        \bottomrule
    \end{tabular}
\end{table}

Des contextes React Native partagent certaines informations entre les écrans, comme l'authentification,
la langue, le thème, les favoris, les filtres et les notifications. Cette approche évite de répéter la même
logique dans plusieurs parties de l'application.

\section{Modèle de données}

Le modèle de données a été construit avec SQLAlchemy \cite{sqlalchemy} à partir des objets réellement
manipulés par l'application. Le tableau suivant ne cherche donc pas à donner une définition générale de ces
entités ; il précise plutôt comment PharmacieConnect organise les comptes, les pharmacies, les gardes, les
médicaments, les sessions, les tentatives de connexion, les recherches et l'audit, comme le détaille le
Tableau~\ref{tab:entites-donnees}.

\begin{longtable}{L{4cm}L{10cm}}
    \caption{Entités principales du modèle de données\\ \small Source : Auteur}\label{tab:entites-donnees}\\
    \toprule
    \textbf{Entité} & \textbf{Description} \\
    \midrule
    \endfirsthead
    \toprule
    \textbf{Entité} & \textbf{Description} \\
    \midrule
    \endhead
    \texttt{Utilisateur} & Compte utilisateur de l'application mobile avec email, mot de passe, statut actif et vérification email. \\
    \texttt{Administrateur} & Compte de gestion avec rôle, statut, informations de profil et périmètre régional éventuel. \\
    \texttt{RefreshToken} & Jeton de rafraîchissement permettant de prolonger une session selon les règles de sécurité. \\
    \texttt{LoginAttempt} & Tentative de connexion utilisée pour suivre les succès et les échecs d'authentification. \\
    \texttt{Pharmacie} & Données d'une pharmacie : nom, adresse, téléphone, gouvernorat, latitude, longitude et référence OSM éventuelle. \\
    \texttt{GardeSchedule} & Planning de garde avec date, pharmacie, horaire, ville, gouvernorat, type de garde et notes. \\
    \texttt{Medicine} & Catalogue de médicaments avec code PCT, nom commercial, prix public, tarif de référence, remboursement, DCI et AP. \\
    \texttt{AuditLog} & Trace des actions importantes avec acteur, cible, statut, adresse IP et détails. \\
    \texttt{SearchEvent} & Événements de recherche utilisés pour alimenter les statistiques et les indicateurs. \\
    \bottomrule
\end{longtable}

\section{Diagramme de classes}

Le diagramme de classes décrit le modèle de données réel de PharmacieConnect, fidèle aux entités persistées
par le backend via SQLAlchemy. Comme le montre la Figure~\ref{fig:class-diagram}, il distingue les comptes
(\texttt{Administrateur}, \texttt{Utilisateur}) et leurs mécanismes de sécurité, les pharmacies et leurs
plannings de garde, les médicaments, ainsi que la traçabilité assurée par les journaux d'audit.

\begin{figure}[H]
    \centering
    \includegraphics[width=0.95\textwidth]{plantuml/class_diagram}
    \caption{Diagramme de classes\\ \small Source : Auteur}
    \label{fig:class-diagram}
\end{figure}

La Figure~\ref{fig:class-diagram} met en évidence les entités centrales du modèle. Les deux types de comptes,
\texttt{Administrateur} (dont le \texttt{role} est typé par \texttt{AdminRoleEnum} et complété d'un
\texttt{region\_scope} pour la délégation régionale) et \texttt{Utilisateur}, partagent les mécanismes de
session : \texttt{RefreshToken} et \texttt{LoginAttempt} y sont rattachés par un lien polymorphe
(\texttt{entity\_type} + \texttt{entity\_id}). L'\texttt{Administrateur} est l'auteur des données qu'il
saisit : les attributs \texttt{created\_by} relient \texttt{Pharmacie}, \texttt{GardeSchedule} et
\texttt{Medicine} au compte qui les a créés, tandis que \texttt{AuditLog} référence son auteur via
\texttt{actor\_id}. La \texttt{Pharmacie} et le \texttt{GardeSchedule} ne sont toutefois pas liés par une
clé étrangère directe : la correspondance s'effectue par le nom de la pharmacie au moment de la requête
(\texttt{pharmacy\_name}). Enfin, \texttt{SearchEvent} conserve les événements de recherche utiles aux
indicateurs. L'assistant de premiers secours n'apparaît pas dans ce schéma, car il fonctionne comme un
service de récupération sans état, sans persistance de conversation. Cette organisation traduit la
séparation des responsabilités : des comptes distincts, des données métier centralisées autour de la
pharmacie et une traçabilité systématique des actions sensibles.

\section{Diagrammes de séquence}

Les diagrammes de séquence retenus décrivent deux parcours critiques du projet. Le premier, présenté dans la
Figure~\ref{fig:sequence-auth}, montre comment une session est créée puis contrôlée à travers les jetons JWT.
Le second, donné par la Figure~\ref{fig:sequence-search}, détaille la recherche d'une pharmacie, depuis la
saisie ou la géolocalisation côté mobile jusqu'au filtrage effectué par le backend. Ces
deux scénarios ont été choisis parce qu'ils concentrent les principaux échanges entre les clients, les
services métier et la base de données.

\begin{figure}[H]
    \centering
    \includegraphics[width=0.95\textwidth]{plantuml/sequence_auth_jwt}
    \caption{Diagramme de séquence : authentification JWT\\ \small Source : Auteur}
    \label{fig:sequence-auth}
\end{figure}

\begin{figure}[H]
    \centering
    \includegraphics[width=0.95\textwidth]{plantuml/sequence_pharmacy_search}
    \caption{Diagramme de séquence : recherche d'une pharmacie\\ \small Source : Auteur}
    \label{fig:sequence-search}
\end{figure}

\section{Diagramme d'activité}

Le diagramme d'activité détaille le traitement d'un import CSV. Ce scénario est important, car il regroupe
plusieurs règles sensibles : contrôle du fichier, validation des colonnes, vérification des lignes,
enregistrement des données valides, rejet des lignes incorrectes et journalisation du résultat
(voir Figure~\ref{fig:activity-import-csv}).

\begin{figure}[H]
    \centering
    \includegraphics[height=0.85\textheight,keepaspectratio]{plantuml/activity_csv_import}
    \caption{Diagramme d'activité : import CSV contrôlé\\ \small Source : Auteur}
    \label{fig:activity-import-csv}
\end{figure}

\section{Diagramme de composants et de déploiement}

Les diagrammes de composants et de déploiement complètent la vision technique. Le premier, représenté dans la
Figure~\ref{fig:component-diagram}, décrit les modules logiciels, tandis que le second, illustré par la
Figure~\ref{fig:deployment-diagram}, présente la répartition logique entre les clients, le serveur backend,
la base de données et la configuration.

\begin{figure}[H]
    \centering
    \includegraphics[width=0.95\textwidth]{plantuml/component_diagram}
    \caption{Diagramme de composants\\ \small Source : Auteur}
    \label{fig:component-diagram}
\end{figure}

\begin{figure}[H]
    \centering
    \includegraphics[width=0.95\textwidth]{plantuml/deployment_diagram}
    \caption{Diagramme de déploiement\\ \small Source : Auteur}
    \label{fig:deployment-diagram}
\end{figure}

\section{Sécurité}
\label{sec:securite}

La sécurité de PharmacieConnect est appliquée côté backend afin que les règles ne dépendent pas uniquement
de l'affichage dans les interfaces. Les routes d'administration exigent un jeton JWT valide, puis vérifient le
rôle du compte connecté. Pour les assistants, le contrôle ne s'arrête pas au rôle : le périmètre régional est
également vérifié avant toute création, modification ou suppression. Cette approche, alignée avec les
recommandations OWASP pour la protection des API~\cite{owasp}, protège particulièrement les imports CSV, les
gardes, les médicaments et les journaux d'audit.

\subsection{Gestion des sessions par jetons JWT}

La gestion des sessions dans PharmacieConnect distingue deux jetons aux durées de vie volontairement
différentes. À la connexion, le backend émet un jeton d'accès valable 15~minutes par défaut, signé avec
HS256, et un jeton de rafraîchissement valable 7~jours. Le premier accompagne chaque requête
protégée dans l'en-tête \texttt{Authorization: Bearer} ; sa courte durée limite les conséquences
d'une interception. Le second est conservé en base dans l'entité \texttt{RefreshToken} et effacé
lors de la déconnexion explicite, ce qui permet de renouveler la session active sans laisser
le jeton d'accès circuler au-delà de sa fenêtre autorisée.

La validation de ces jetons est prise en charge exclusivement par le backend à chaque appel
protégé : déléguer ce contrôle à l'interface cliente serait insuffisant, car tout appel direct à
l'API contournerait alors les vérifications \cite{jwt}.

\subsection{Contrôle d'accès basé sur les rôles (\gls{rbac})}

PharmacieConnect organise les accès administratifs selon trois rôles : \texttt{super\_admin},
\texttt{admin} et \texttt{assistant}. Chaque endpoint protégé évalue le rôle du demandeur avant d'exécuter
l'opération. Pour le rôle \texttt{assistant}, un deuxième filtre s'applique : l'attribut
\texttt{region\_scope} limite les actions au périmètre régional autorisé. Toute tentative d'opération hors
de ce périmètre est rejetée même si le jeton d'accès est valide. Cette combinaison rôle\,/\,périmètre répond
directement au besoin de délégation régionale : plusieurs assistants peuvent coexister sans modifier les
données qui ne relèvent pas de leur zone.

\subsection{Hachage des mots de passe}

Le projet a retenu bcrypt via Passlib pour protéger les mots de passe des administrateurs et des
utilisateurs mobiles. Chaque empreinte est produite avec un sel unique généré à l'enregistrement :
deux comptes utilisant le même mot de passe donnent des empreintes différentes, ce qui empêche
toute comparaison directe entre entrées. Le facteur de coût est réglé pour que la vérification
lors de la connexion reste imperceptible pour l'utilisateur, tout en rendant une attaque par
volume prohibitivement longue. La comparaison à la connexion s'opère uniquement sur les empreintes
stockées ; le mot de passe en clair n'est jamais conservé en base \cite{passlib_bcrypt}.

\subsection{Validation des entrées avec Pydantic}

Avant qu'une donnée entrante n'atteigne la couche service, un schéma Pydantic la valide
intégralement. Pour PharmacieConnect, cela concerne notamment les formats d'email des comptes,
les plages de coordonnées des pharmacies (latitude bornée entre $-90$ et $90$, longitude entre
$-180$ et $180$), les dates des plannings de garde, les codes PCT des médicaments et la taille
maximale des fichiers CSV importés. Regrouper ces contrôles dans les schémas évite de les
répliquer dans chaque route et intercepte les données malformées avant tout accès à la base \cite{pydantic}.

\subsection{Politique \gls{cors} et hôtes autorisés}

Le backend déclare explicitement les origines autorisées via les variables d'environnement de configuration
CORS. Toute origine non prévue par la configuration ne doit pas accéder aux routes depuis un navigateur.
En complément, \texttt{TrustedHostMiddleware} inspecte l'en-tête d'hôte de chaque requête et bloque celles
dont l'hôte n'est pas inscrit dans \texttt{TRUSTED\_HOSTS}. Ces paramètres sont externalisés afin d'adapter
la configuration entre le développement local et un déploiement en production.

\subsection{Traçabilité et journalisation des actions}

Deux mécanismes de traçabilité coexistent dans PharmacieConnect. La table \texttt{LoginAttempt} enregistre
les tentatives de connexion avec leur statut, l'adresse IP source et l'horodatage. En parallèle, le service
\texttt{AuditService} journalise les opérations administratives significatives : ajout, mise à jour et
retrait d'entités, imports CSV et changements de droits. Chaque entrée d'audit identifie l'acteur,
décrit l'opération, référence la cible et indique le résultat. Cette granularité permet de reconstituer
l'historique des modifications apportées aux données de la plateforme.

\subsection{Sécurité du déploiement}

Pour un déploiement en production, les mesures complémentaires recommandées sont :

\begin{itemize}
    \item forcer \textbf{HTTPS} sur toutes les communications et rediriger le trafic HTTP ;
    \item utiliser une \texttt{SECRET\_KEY} aléatoire d'au moins 256 bits, distincte entre
    les environnements de développement et de production ;
    \item remplacer SQLite par \textbf{PostgreSQL} pour une gestion robuste des connexions
    concurrentes ;
    \item activer le cache \textbf{Redis} pour réduire la pression sur la base de données et
    limiter les requêtes répétitives.
\end{itemize}

\section{Justification des choix techniques}

Le Tableau~\ref{tab:choix-techniques} justifie les principaux choix techniques retenus pour le projet.

\begin{longtable}{L{4cm}L{10cm}}
    \caption{Justification des choix techniques\\ \small Source : Auteur}\label{tab:choix-techniques}\\
    \toprule
    \textbf{Choix} & \textbf{Justification} \\
    \midrule
    \endfirsthead
    \toprule
    \textbf{Choix} & \textbf{Justification} \\
    \midrule
    \endhead
    FastAPI & Retenu pour concentrer les endpoints publics et protégés dans un backend testable, avec validation des requêtes et documentation API utile pendant les essais \cite{fastapi}. \\
    React et Vite & Retenus pour structurer le tableau de bord par pages fonctionnelles : pharmacies, gardes, médicaments, audit et paramètres \cite{react,vite}. \\
    Expo et React Native & Retenus pour réaliser les écrans mobiles liés au parcours citoyen : recherche, carte, calendrier, médicaments, profil et assistant \cite{expo,reactnative}. \\
    SQLAlchemy & Retenu pour représenter les entités métier du projet et garder l'accès aux tables dans les services backend \cite{sqlalchemy}. \\
    SQLite et PostgreSQL & SQLite facilite la validation locale, tandis que PostgreSQL reste plus adapté à un déploiement avec connexions concurrentes \cite{postgresql}. \\
    JWT & Retenu pour séparer un jeton d'accès court et un jeton de rafraîchissement, afin de protéger les routes administratives et les comptes utilisateurs \cite{jwt}. \\
    CSV & Retenu comme format d'import pragmatique pour les pharmacies, les gardes et les médicaments, avec validation des colonnes et retour des erreurs. \\
    \bottomrule
\end{longtable}

\section{Conclusion}

L'ensemble de cette conception offre une base cohérente pour passer à la réalisation de PharmacieConnect.
Les règles sensibles demeurent côté backend, chaque interface se consacre à son propre parcours et les
données métier s'organisent autour des pharmacies, des gardes, des médicaments et de l'audit. Une telle
séparation facilite la maîtrise des évolutions futures : ajout d'une règle d'import, création d'un rôle
administratif ou perfectionnement de l'assistant de premiers secours. Le chapitre qui suit fait le pont
entre cette conception, la réalisation technique et la validation de la solution.
